\subsection{Descent through the actual Galois quotient}
\label{sec:support-finite-end}

Retain an original \texttt{GalLink} $k$ from
$(\eta,\Psi,\sigma,g_1,B,M,L)$ to a frame $F$ satisfying all its
original \texttt{Ok} hypotheses, where $\ell\mid M$ and $\sigma^2=1$.
The field $L$ is the specified number field; it is CM when supplied
by \texttt{HasFrame}.  Write $U^+=\operatorname{Gal}(\overline{\mathbb Q}/L^+)$
for the original subgroup fixing its maximal real subfield.
The exact \texttt{GalLink} includes a continuous surjection
$p:U^+\to\Gamma$ and
$\operatorname{Gal}(\overline{\mathbb Q}/L)=U^+\cap\ker\eta$.
All point and quotient conditions here are those of
\cite{support-finite-end-source}.

\begin{theorem}[The original normalized lift and finite-end transfer]
\label{thm:support-finite-end}
Every original \texttt{OrdPt} $r:G_{\mathbb Q}\to GG_3(A)$ over
a Noetherian local domain $A$ with local $O$-map gives an original
normalized lift $\rho$ over the same $A$, with
\[
 \rho(p_0(s))=r(s)_{\rm mat}\quad(s\in G_L),\qquad
 \psi_{v,i}(g)=\Psi(j_v(g))^{1-i}\quad(g\in I_v),            \tag{FE1}
\]
where the rightmost values use the prescribed map $O\to A$.
If $\operatorname{FinPt}(O,F,3)$ holds, then the literal
\texttt{FinEnd} holds with
\[
 U=G_L,\qquad [G_{\mathbb Q}:U]=[L:\mathbb Q].              \tag{FE2}
\]
This one $U$ works for every original complete DVR $A$ killing
$\mathfrak m_O$ and every original ordinary point $r$ over it.
The theorem retains the actual \texttt{GalLink} as an input.
\end{theorem}

\begin{proof}
The original kernel of $p$ is the closed normal closure in $U^+$
of inertia elements away from $\ell$ with $\eta=\overline\Psi=1$.
The \texttt{GalLink.teich} condition gives $\Psi=1$ there, and
\texttt{OrdPt.away} makes $r=1$ on these generators and their
normal closure.  The matrix-one locus is closed: it is the
intersection of the closed fibres where the matrix is one modulo
$\mathfrak m_A^N$, and Krull intersection identifies that
intersection.  The multiplier on $\ker p$ is $\eta=1$, so
$\ker p\subset\ker(r|_{U^+})$ for the whole semidirect product.
Thus $r$ descends to $r':\Gamma\to GG_3(A)$ with
$r'(p(s))=r(s)$ and multiplier $\delta$.  It remains adically
continuous, since the continuous surjection from compact $U^+$
to Hausdorff $\Gamma$ is a quotient map and every matrix reduction
is locally constant after pullback.  The matrix part on
$G_0=\ker\delta$ is the required representation $\rho$.

The exact marked identities $p(\sigma)=c$, $p(g_1)=g_0$ and
$r(\sigma)_{\rm mat}=1$, together with its multiplier $-1$, give
\[
 \rho(\operatorname{star}g)=J\rho(g)^tJ.                    \tag{FE3}
\]
The two original equations $z_{12}=z_{21}=0$ at $g_1$ make every
off-diagonal entry of $\rho(g_0)-\rho(\operatorname{star}g_0)$
zero: in $X-JX^tJ$ the two corner off-diagonal entries are
identically zero, and the remaining four occur as two opposite
pairs.  The residual matrix and diagonal characters descend
unchanged by the exact identity
$\Psi_F(p(s))=\overline\Psi(s)$.

For each $v$, the original place correspondence supplies
$\gamma$ with $\eta(\gamma)=1$ and $\operatorname{pl}(v)=\gamma B$.
Then $d(g)=\gamma^{-1}j_v(g)\gamma$ lies in $D_B$, and lies in
$I_B$ when $g\in I_v$.  Its multiplier is one by the exact
local map to $G_0$.  The original \texttt{CondT} gives
$h\in\operatorname{GL}_3(A)$ with $h\equiv1\pmod{\mathfrak m_A}$
such that
\[
 T(g)=h^{-1}r(d(g))_{\rm mat}h
\]
is upper triangular and has its specified diagonal on inertia.
Set $\psi_{v,i}(g)=T(g)_{ii}$.  These are continuous unit-valued
characters, with the exact inertia values in (FE1), because
$\Psi$ is invariant under conjugation.  Moreover
$r(j_v(g))_{\rm mat}=PT(g)P^{-1}$ for the integral invertible
matrix $P=r(\gamma)_{\rm mat}h$.

For arbitrary $g_0,g_1,g_2$ in this local group,
\[
 (T(g_0)-T(g_0)_{00}I)(T(g_1)-T(g_1)_{11}I)
                         (T(g_2)-T(g_2)_{22}I)=0.          \tag{FE4}
\]
Starting on the right, the three factors take the full space to
the second flag step, to the first step, and then to zero.
Conjugation gives the literal ordered local product for $\rho$;
(FE3) gives its polarized companion with the correct reversal.
The congruence $h\equiv1$ preserves the prescribed residual
characters.  This proves the complete normalized lift and (FE1)
without changing the coefficient ring or lattice.

For the final assertion, choose $U=G_L$ before choosing $A$ or $r$.
If $A$ is an original DVR killing $\mathfrak m_O$, its $O$-map
factors through $k(O)$, giving the exact scalar tower for
\texttt{FinPt}.  For $g\in I_v\cap\ker\Psi_F$, (FE1) and the
exact character correspondence give $\psi_{v,i}(g)=1$.
Apply \texttt{FinPt} and then the first identity of (FE1): for
every $s\in U$, the characteristic polynomial of $r(s)_{\rm mat}$
equals that of the original residual $r_0(\eta,\Psi)(s)$ mapped
to $A$.  This is precisely \texttt{FinEnd}, uniformly in every
tested $A$ and $r$, with index as in (FE2).
\end{proof}

No bound on the degree of the auxiliary field is inferred from
its finiteness.  The original auxiliary frame, including its growing
local degrees, class-rank bound, two strict cohomology classes and seed,
is constructed in Theorem~\ref{thm:support-auxiliary-frame}.
